\chapter{Existence Theory for the Ricci Flow}

\section*{5.1 \quad Ricci flow is not parabolic}

\begin{remark}
\label{topping-ricci-flow-linearisation-not-parabolic}
\source{topping:page-0060,topping:page-0061}
\uses{topping:variation-ricci, topping:lie-derivative-symmetric-gradient}
\dcref{ch5:1.1}
\group{topping-ricci-flow-existence}

On $E=\Sym^2T^*\M$, write the Ricci-flow operator as
\[
Q(g)=-2\Ric(g).
\]
Its linearisation is
\[
\frac{\partial h}{\partial t}=Lh=[DQ(g)]h
 =\Delta_{\mathcal L}h+\mathcal L_{(\dG(h))^\#}g.
\]
For $(x,\xi)\in T^*\M$,
\[
\sigma(L)(x,\xi)h
 =|\xi|^2h-\xi\otimes h(\xi^\sharp,\cdot)
 -h(\xi^\sharp,\cdot)\otimes\xi+(\xi\otimes\xi)\tr h.
\]
This symbol vanishes when $h=\xi\otimes\xi$, so the linearised Ricci-flow
equation is not parabolic.
\end{remark}

\begin{proof}
\sketch The Lichnerowicz Laplacian has the same principal symbol as the
rough Laplacian, namely $\sigma(\Delta_{\mathcal L})(x,\xi)h=|\xi|^2h$.
Using $\dG(h)=\delta h+\frac12d(\tr h)$ and
\[
\mathcal L_{\omega^\#}g(X,W)=\nabla\omega(X,W)+\nabla\omega(W,X),
\]
the correction term contributes
\[
\sigma\!\Bigl(h\mapsto\mathcal L_{(\dG(h))^\#}g\Bigr)(x,\xi)h
 =-\xi\otimes h(\xi^\sharp,\cdot)
 -h(\xi^\sharp,\cdot)\otimes\xi+(\xi\otimes\xi)\tr h.
\]
Adding the two symbols gives the displayed formula, whose value at
$h=\xi\otimes\xi$ is zero.
\end{proof}

\section*{5.2 \quad Short-time existence and uniqueness: the DeTurck trick}

Fix a smooth positive-definite tensor
$T\in\Gamma(\Sym^2T^*\M)$ and define
\[
P(g)=-2\Ric(g)+\mathcal L_{T^{-1}\dG(T)}^{\#}g.
\tag{5.2.1}
\]
The second-order part of $DP(g)$ is $\Delta$, so $P$ is parabolic at every
metric.

\begin{theorem}[Short time existence]
\label{topping-ricci-flow-short-time-existence}
\source{topping:page-0062,topping:page-0063}
\uses{topping:variation-ricci, topping:lie-derivative-symmetric-gradient}
\dcref{ch5:2.1}
\group{topping-ricci-flow-existence}

For every smooth metric $g_0$ on a closed manifold $\M$, there are
$\epsilon>0$ and a smooth metric family $g(t)$ on $[0,\epsilon]$ satisfying
\[
\begin{cases}
\dfrac{\partial g}{\partial t}=-2\Ric(g) & t\in[0,\epsilon],\\
g(0)=g_0.
\end{cases}
\tag{5.2.2}
\]
\end{theorem}

\begin{proof}
\sketch Solve the parabolic DeTurck equation $\partial_tg=P(g)$ with
initial metric $g_0$. The flow of
\[
X=-(T^{-1}\dG(T))^\#
\]
produces diffeomorphisms $\psi_t$ with $\psi_0=\id$. Pulling the DeTurck
solution back by $\psi_t$ removes the Lie-derivative term and gives the
required Ricci flow.
\end{proof}

\begin{theorem}[Uniqueness of solutions (forwards and backwards)]
\label{topping-ricci-flow-uniqueness}
\source{topping:page-0064}
\dcref{ch5:2.2}
\group{topping-ricci-flow-existence}
If $g_1(t)$ and $g_2(t)$ are Ricci flows on a closed manifold $\M$ over
$[0,\epsilon]$, with $\epsilon>0$, and $g_1(s)=g_2(s)$ for some
$s\in[0,\epsilon]$, then
\[
g_1(t)=g_2(t)\qquad\text{for every }t\in[0,\epsilon].
\]
\end{theorem}

\begin{proof}
\sketch Put each flow into DeTurck gauge using harmonic map flow into the
fixed background metric $T$, starting from the identity map. The gauged flows
satisfy the same parabolic equation with the same data at the specified time;
parabolic uniqueness in the forward and backward directions makes the gauges,
and hence the original flows after pushing forward, agree on the whole
interval.
\end{proof}

Thus a Ricci flow with initial metric $g_0$ is defined on a maximal interval
$[0,T]$ (or $[0,T)$ when the endpoint is singular).

\section*{5.3 \quad Curvature blow-up at finite-time singularities}

\begin{lemma}[Metric equivalence]
\label{topping-ricci-flow-metric-equivalence}
\source{topping:page-0065}
\dcref{ch5:3.2}
\group{topping-ricci-flow-existence}
If $g(t)$ is a Ricci flow on $[0,s]$ and $|\Ric|\leq M$ on
$\M\times[0,s]$, then for every $t\in[0,s]$,
\[
e^{-2Mt}g(0)\leq g(t)\leq e^{2Mt}g(0).
\]
\end{lemma}

\begin{proof}
\sketch For every nonzero tangent vector $X$,
\[
\frac{\partial}{\partial t}g(X,X)=-2\Ric(X,X),
\qquad
\left|\frac{\partial}{\partial t}\ln g(X,X)\right|\leq 2M.
\]
Integration gives
\[
\left|\ln\frac{g(t)(X,X)}{g(0)(X,X)}\right|\leq 2Mt,
\]
which is equivalent to the two-sided tensor estimate.
\end{proof}

\begin{theorem}[Curvature blows up at a singularity]
\label{topping-ricci-flow-curvature-blows-up}
\source{topping:page-0064,topping:page-0065,topping:page-0066,topping:page-0067}
\uses{topping-maxpr-curvature-norm-bound, topping-maxpr-higher-derivatives, topping-ricci-flow-metric-equivalence}
\dcref{ch5:3.1}
\group{topping-ricci-flow-existence}

If $\M$ is closed and $g(t)$ is a Ricci flow on a maximal interval
$[0,T)$ with $T<\infty$, then
\[
\sup_{\M}|\Rm|(\cdot,t)\longrightarrow\infty
\qquad\text{as }t\uparrow T.
\]
\end{theorem}

\begin{proof}
\sketch If $|\Rm|$ were bounded on $[0,T)$, Theorem~\ref{topping-maxpr-curvature-norm-bound} would give a
uniform bound $|\Rm|\leq M$ on $\M\times[0,T)$. The Ricci bound and the
metric-equivalence lemma would then imply that $g(t)$ converges continuously
to a positive-definite tensor $g(T)$ as $t\uparrow T$.

On $[\frac12T,T)$, Corollary~\ref{topping-maxpr-higher-derivatives} bounds all quantities
$\partial_t^\ell\nabla^k\Rm$. Together with the Ricci-flow equation and the
commutator formula for $\partial_t$ and $\nabla$, these bounds control all
mixed space-time derivatives of the coordinate coefficients $g_{ij}$. Hence
$g(t)$ extends smoothly to $T$. Applying the short-time existence theorem to
$g(T)$ extends the flow beyond $T$, contradicting maximality.
\end{proof}
